\section{Externí datové a informační zdroje. Sekundární zdroje dat a informací, role competitive intelligence, přístup k externím zdrojům}

\subsection*{Smysl externích zdrojů}
Externí datový nebo informační zdroj je zdroj, který nevzniká primárně uvnitř analyzované organizace.
Může jít o veřejnou statistickou databázi, registr veřejné správy, otevřená data, placenou databázi,
oborový report, data z API, webové stránky, dokumenty konkurentů, vědecké publikace nebo mediální
monitoring.

Externí zdroje se v datové analytice používají hlavně k tomu, aby interní data dostala širší kontext.
Samotné interní tržby například neřeknou, zda firma roste díky vlastní výkonnosti, nebo jen kopíruje
růst celého trhu. Po doplnění externích dat lze porovnávat výkon firmy s trhem, hodnotit regionální
potenciál, odhadovat poptávku, sledovat makroekonomické vlivy, obohacovat zákaznická data nebo
vytvářet varovné signály o změnách v okolí podniku.

\subsection*{Data, informace a zdroje}
Pro zkoušku je dobré rozlišit tři roviny:
\begin{itemize}
  \item \term{data} - jednotlivé hodnoty nebo záznamy, například počet obyvatel obce, IČO firmy,
  kurz měny nebo cena produktu;
  \item \term{informace} - data zasazená do významu, například pokles počtu zákazníků v regionu
  s klesající kupní silou;
  \item \term{znalost pro rozhodnutí} - interpretace informace ve vztahu k cíli, například zda má
  firma otevřít novou pobočku, změnit cenu nebo upravit segmentaci zákazníků.
\end{itemize}

Externí zdroj proto nemá hodnotu jen tím, že existuje. Hodnotu má tehdy, když je důvěryhodný,
srozumitelný, legálně použitelný, technicky dostupný a relevantní pro konkrétní analytickou otázku.

\subsection*{Primární a sekundární zdroje}
\term{Primární zdroj} přináší původní data nebo původní informaci. V datové analytice to mohou být
přímo naměřená data, odpovědi respondentů, experimentální data, provozní logy, interní transakce,
senzorová měření nebo originální dokumenty organizace. U externích zdrojů mohou být primárním zdrojem
například úřední registry, účetní závěrky ve sbírce listin, oficiální statistické časové řady,
meteorologická měření nebo data publikovaná institucí, která je sama vytváří.

\term{Sekundární zdroj} původní data nevytváří, ale shrnuje, interpretuje, agreguje nebo odkazuje
na jiné zdroje. Patří sem slovníky, encyklopedie, lexikony, přehledové články, analytické reporty,
tržní studie, oborové databáze, agregátory cen, katalogy datových sad nebo rešerše.

Rozdíl není vždy absolutní. Datový portál může být sekundárním zdrojem metadat, protože katalogizuje
datové sady jiných institucí, ale samotná datová sada zveřejněná přes tento portál může být primárním
zdrojem pro analýzu. U zkoušky je dobré vysvětlit nejen typ zdroje, ale i to, kdo data vytvořil,
kdo je publikuje a zda jde o originální měření, nebo o odvozenou interpretaci.

\subsection*{Typy externích zdrojů}
\begin{center}
\small
\begin{tabular}{@{}>{\raggedright\arraybackslash}p{0.24\textwidth}
                >{\raggedright\arraybackslash}p{0.35\textwidth}
                >{\raggedright\arraybackslash}p{0.30\textwidth}@{}}
\hline
\term{Typ zdroje} & \term{Příklady} & \term{Analytické využití} \\
\hline
Oficiální statistiky & ČSÚ, Eurostat, OECD, World Bank, ČNB ARAD &
makroekonomické proměnné, regiony, demografie, benchmarky \\
Veřejné registry & ARES, veřejný rejstřík, registr smluv, RÚIAN, insolvenční rejstřík &
identifikace firem, vlastnické vazby, sídla, veřejné zakázky, riziko partnera \\
Open data portály & data.gov.cz, data.europa.eu, městské a krajské portály &
strojově čitelné datové sady, geodata, doprava, životní prostředí \\
Komerční databáze & tržní databáze, ratingové databáze, kreditní informace, panelová data &
tržní podíly, bonita, konkurence, cílení kampaní \\
Web a dokumenty & weby firem, ceníky, e-shopy, výroční zprávy, tiskové zprávy, nabídky práce &
competitive intelligence, cenový monitoring, signály růstu nebo změny strategie \\
Sociální a mediální data & sociální sítě, zprávy, recenze, diskuse, mediální monitoring &
sentiment, reputace, témata zákazníků, krizové signály \\
Geografická a senzorová data & mapy, družicová data, počasí, dopravní data, IoT zdroje &
lokalizační analýzy, logistika, predikce poptávky, rizika počasí \\
\hline
\end{tabular}
\end{center}

\subsection*{Příklady veřejných zdrojů pro ČR a EU}
U státnic je užitečné umět uvést konkrétní zdroje a říct, k čemu se hodí:
\begin{itemize}
  \item \term{ČSÚ} - oficiální statistiky o obyvatelstvu, ekonomice, regionech, cenách, trhu práce
  a podnikatelském sektoru; část dat je dostupná přes veřejnou databázi a otevřená data.
  \item \term{Eurostat} - evropské statistiky pro srovnání států a regionů; data lze stahovat přes
  Data Browser i API.
  \item \term{data.gov.cz} - Národní katalog otevřených dat, který pomáhá najít datové sady veřejné
  správy a jejich metadata.
  \item \term{data.europa.eu} - oficiální evropský portál dat; sjednocuje metadata z institucí EU
  i členských států a obsahuje také datové sady s vysokou hodnotou.
  \item \term{ARES} - administrativní registr ekonomických subjektů; umožňuje vyhledávat údaje
  o ekonomických subjektech registrovaných v ČR a nabízí veřejné API.
  \item \term{Veřejný rejstřík a sbírka listin} - zákonem stanovené údaje o právnických osobách,
  výpisy, účetní závěrky, výroční zprávy a další listiny.
  \item \term{ČNB ARAD} - veřejná databáze časových řad z oblasti měnové, finanční a ekonomické
  statistiky; pro automatizované použití existuje ARAD-REST-API.
  \item \term{RÚIAN a geodata ČÚZK} - adresy, územní prvky, parcely a další prostorová data.
  \item \term{Registr smluv, NEN a Věstník veřejných zakázek} - veřejné smlouvy a zakázky, vhodné
  pro analýzu veřejných výdajů, dodavatelů a trhu B2G.
\end{itemize}

\subsection*{Otevřená data}
\term{Otevřená data} jsou data zveřejněná tak, aby je bylo možné snadno najít, stáhnout, strojově
zpracovat a znovu použít. Důležité vlastnosti otevřených dat jsou:
\begin{itemize}
  \item dostupnost na webu bez zbytečných překážek;
  \item strojově čitelný a otevřený formát, například CSV, JSON, XML nebo RDF;
  \item jasné podmínky užití nebo licence;
  \item evidence v datovém katalogu;
  \item metadata, dokumentace a kontakt na správce dat;
  \item možnost hromadného stažení nebo přístupu přes API, pokud to povaha dat vyžaduje.
\end{itemize}

Je potřeba odlišit otevřená data od běžné webové stránky. Vyhledávací formulář, PDF určené pro čtení
člověkem nebo mapa bez možnosti získat datovou vrstvu nejsou plnohodnotná otevřená data. Mohou být
užitečné jako informační zdroj, ale pro analytickou pipeline jsou hůře automatizovatelné.

V EU se zvlášť zdůrazňují \term{datové sady s vysokou hodnotou}. Patří do oblastí jako geodata,
pozorování Země a životní prostředí, meteorologie, statistika, společnosti a vlastnictví společností
a mobilita. Mají být dostupné ve strojově čitelné podobě, přes API a hromadné stažení a s otevřenou
licencí.

\subsection*{Hodnocení důvěryhodnosti a kvality}
Před použitím externího zdroje nestačí ověřit, že data jdou stáhnout. Je nutné posoudit, zda zdroj
odpovídá účelu analýzy.

Otázky k důvěryhodnosti zdroje:
\begin{itemize}
  \item kdo je autorem nebo správcem zdroje;
  \item zda jde o oficiální instituci, odbornou organizaci, univerzitu, komerčního poskytovatele
  nebo neověřený web;
  \item zda je uvedena metodika sběru, definice proměnných, časový rozsah a periodicita aktualizace;
  \item zda zdroj uvádí metadata, dokumentaci, licenci a kontaktní osobu;
  \item zda lze data ověřit proti jinému zdroji;
  \item zda nejsou data jednostranná, marketingově upravená nebo vytržená z kontextu.
\end{itemize}

Otázky k datové kvalitě:
\begin{itemize}
  \item \term{úplnost} - nechybí klíčové záznamy, období, regiony nebo proměnné;
  \item \term{přesnost} - odpovídají hodnoty skutečnosti a metodice měření;
  \item \term{konzistence} - nejsou v datech rozpory mezi tabulkami, číselníky a jednotkami;
  \item \term{aktuálnost} - jsou data dostatečně čerstvá pro rozhodnutí;
  \item \term{validita} - měří proměnná opravdu to, co má měřit;
  \item \term{unikátnost} - nevznikají duplicity jednotek;
  \item \term{integrita} - sedí vazby mezi identifikátory, například IČO, kód obce nebo NUTS region.
\end{itemize}

Častý problém externích dat je rozdílná granularita. Interní data mohou být denní po zákaznících,
zatímco externí statistika je měsíční po regionech. Analytik pak musí rozhodnout, zda bude data
agregovat, spojovat přes geografický klíč, interpolovat, nebo je použije jen jako kontext.

\subsection*{Právní a licenční omezení}
Externí data nejsou automaticky volně použitelná. Vždy je potřeba ověřit:
\begin{itemize}
  \item \term{licenci} - zda je dovoleno data kopírovat, upravovat, kombinovat, sdílet a použít
  komerčně;
  \item \term{autorská a databázová práva} - zejména u databází, reportů, textů, obrázků a map;
  \item \term{smluvní podmínky} - API může mít limity, zákaz redistribuce nebo povinnost uvést zdroj;
  \item \term{ochranu osobních údajů} - pokud data obsahují nebo mohou odhalit fyzickou osobu;
  \item \term{obchodní tajemství a důvěrnost} - zejména u nakoupených, partnerských nebo neveřejných dat;
  \item \term{pravidla web scrapingu} - podmínky webu, soubor \texttt{robots.txt}, zákaz obcházení
  technických ochran a přiměřená zátěž serveru.
\end{itemize}

U otevřených dat bývá licence nastavena tak, aby umožnila opakované použití. To však neznamená, že
odpadá odpovědnost za správnou interpretaci. I otevřený dataset může mít omezenou přesnost, neúplná
metadata nebo metodické změny v čase.

\subsection*{Přístup k externím zdrojům}
Přístup k externím zdrojům může mít různé podoby:
\begin{itemize}
  \item \term{ruční stažení souboru} - jednoduché pro jednorázovou analýzu, ale špatně reprodukovatelné;
  \item \term{hromadné stažení} - vhodné pro pravidelné načítání celých datasetů;
  \item \term{API} - strukturované rozhraní pro automatizovaný přístup, často s autentizací,
  stránkováním a limity;
  \item \term{datový katalog} - zdroj metadat, odkazů, licencí a dokumentace;
  \item \term{přímý databázový přístup} - typický u partnerských nebo placených zdrojů;
  \item \term{web scraping} - automatizované získávání dat z HTML stránky, pokud neexistuje vhodné
  API a pokud to dovolují právní, technická a etická omezení;
  \item \term{nákup dat} - komerční licence k datům, často s lepší dokumentací, podporou a SLA.
\end{itemize}

Pokud existuje API, je obvykle lepší než scraping. API je stabilnější, má jasnější strukturu, může
vracet metadata a bývá navrženo pro strojové zpracování. Scraping je křehčí, protože změna HTML
šablony může rozbít extrakční pravidla.

\subsection*{Napojení externích zdrojů do datové pipeline}
Externí data je vhodné zapojovat stejně disciplinovaně jako interní zdroje. Typický tok:
\begin{enumerate}
  \item výběr zdroje a ověření licence, metodiky a periodicity aktualizace;
  \item stažení souboru, hromadný export nebo volání API do staging vrstvy;
  \item validace kvality, schématu, počtu řádků, identifikátorů a časového rozsahu;
  \item transformace a napojení na interní data;
  \item využití v reportingu, modelu, dashboardu nebo rozhodovacím procesu.
\end{enumerate}

Praktické zásady:
\begin{itemize}
  \item ukládat surovou kopii dat, datum stažení, URL, verzi a licenci;
  \item oddělit technické načtení od analytické transformace;
  \item kontrolovat schéma, datové typy, počet řádků, duplicity a rozsah hodnot;
  \item používat stabilní identifikátory, například IČO, kód obce, NUTS, ISO kód státu nebo časovou
  dimenzi;
  \item monitorovat změny API, výpadky, limity dotazů a metodické změny;
  \item dokumentovat, které externí proměnné byly použity a jak ovlivňují výsledek.
\end{itemize}

Například při obohacení zákaznických dat lze interní adresu napojit na RÚIAN nebo kód obce, k tomu
přidat demografii z ČSÚ, regionální ukazatele z Eurostatu a vzdálenost k pobočce z geodat. Výsledek
může sloužit k segmentaci, predikci poptávky nebo plánování pobočkové sítě.

\subsection*{Competitive intelligence}
\term{Competitive intelligence} je systematické získávání, třídění, analýza a předávání informací
o vnějším prostředí organizace. Nejde jen o sledování přímých konkurentů. Patří sem také zákazníci,
dodavatelé, substituty, technologie, regulace, makroekonomické trendy, nové obchodní modely a změny
v chování trhu.

Cílem competitive intelligence není vytvořit co největší archiv dokumentů. Cílem je dodat včasnou,
ověřenou a použitelnou informaci pro rozhodnutí. Typické otázky jsou:
\begin{itemize}
  \item vstupuje na trh nový konkurent;
  \item mění konkurence ceny, produktové balíčky nebo distribuční kanály;
  \item roste poptávka v určitém segmentu nebo regionu;
  \item objevuje se technologie, která může změnit náklady nebo kvalitu služby;
  \item hrozí regulatorní změna, která ovlivní byznys model;
  \item je konkrétní dodavatel nebo odběratel finančně rizikový.
\end{itemize}

Competitive intelligence je legální a etická disciplína, pokud pracuje s veřejně dostupnými,
legálně získanými nebo smluvně dostupnými informacemi. Není to průmyslová špionáž, hacking,
obcházení přístupových práv ani získávání důvěrných informací pod falešnou identitou.

\subsection*{Proces competitive intelligence}
Proces lze popsat jako cyklus:
\begin{enumerate}
  \item \term{Plánování a zadání} - vyjasnit rozhodnutí, které má být podpořeno. Formulují se
  \term{KIQ - Key Intelligence Questions}, tedy klíčové zpravodajské otázky.
  \item \term{Mapování zdrojů} - určit, které veřejné registry, databáze, weby, dokumenty, reporty,
  rozhovory nebo placené zdroje mohou odpověď poskytnout.
  \item \term{Sběr dat a informací} - získat data přes API, katalogy, dokumenty, monitoring webu,
  rešerše nebo legální primární výzkum.
  \item \term{Vyhodnocení zdrojů} - posoudit důvěryhodnost, aktuálnost, zaujatost, metodiku a právní
  použitelnost.
  \item \term{Analýza} - převést data na závěry. Používá se například benchmarking, trendová analýza,
  SWOT, PESTLE, Porterův model pěti sil, scénáře nebo early warning indikátory.
  \item \term{Distribuce} - předat závěry ve formě reportu, dashboardu, alertu, briefing note nebo
  doporučení pro management.
  \item \term{Zpětná vazba} - ověřit, zda výstup pomohl rozhodnutí, a zpřesnit další sběr.
\end{enumerate}

Dobrá competitive intelligence odděluje fakta od interpretací. Fakt je například zveřejněná účetní
závěrka, otevření nové pobočky nebo změna ceníku. Interpretace je závěr, že konkurent posiluje určitý
segment, má problém s marží nebo se připravuje na expanzi. U interpretace je nutné uvést, o jaké
důkazy se opírá a jaké existují alternativní výklady.

\subsection*{Příklady analytických úloh}
\term{Cenový monitoring} sleduje veřejné ceny konkurentů, slevy, dostupnost a produktové balíčky.
Data mohou pocházet z webů, e-shopů, ceníků nebo API agregátorů. Rizikem je, že veřejná cena nemusí
odpovídat skutečně realizované ceně po individuálních slevách.

\term{Tržní a regionální analýza} kombinuje interní prodeje s externími daty o populaci, příjmech,
zaměstnanosti, dopravní dostupnosti a konkurenci v území. Používá se pro volbu lokality, obchodní
plán nebo rozdělení prodejních cílů.

\term{Riziko obchodního partnera} využívá registry firem, insolvenční rejstřík, účetní závěrky,
ARES, sankční seznamy, platební historii a mediální monitoring. Výstupem může být skóre rizika,
varování nebo pravidlo pro schvalování obchodního vztahu.

\term{Technologický scouting} sleduje patenty, vědecké články, startupy, pracovní inzeráty, konference
a investice konkurentů. Cílem je odhalit technologické trendy dříve, než se projeví ve finančních
výsledcích.

\subsection*{Rizika a dobré praktiky}
\begin{itemize}
  \item \term{Falešná přesnost:} externí dataset může vypadat přesně, ale být založený na odhadech,
  vzorku nebo změněné metodice.
  \item \term{Nesoulad definic:} tržba, zákazník, domácnost, firma nebo region mohou být v různých
  zdrojích definovány odlišně.
  \item \term{Zpoždění dat:} oficiální statistiky bývají kvalitní, ale často vycházejí se zpožděním.
  \item \term{Selektivní výběr zdrojů:} výsledek lze zkreslit tím, že analytik vybírá jen zdroje
  potvrzující předem zvolený závěr.
  \item \term{Křehká automatizace:} scraping nebo nestabilní API se může bez varování změnit.
  \item \term{Právní riziko:} licence nebo smluvní podmínky mohou zakazovat redistribuci, automatizované
  stahování nebo komerční použití.
  \item \term{Chybné spojování dat:} špatné mapování identifikátorů, firem, adres nebo regionů může
  vytvořit přesvědčivě vypadající, ale chybný výsledek.
\end{itemize}

\pozor
U zkoušky je dobré neplést si externí zdroj s automaticky kvalitním zdrojem. Externí data mohou být
oficiální, otevřená i strojově čitelná, a přesto nemusí být vhodná pro danou úlohu. V odpovědi je
vhodné vždy zmínit původ dat, metodiku, licenci, kvalitu, technický přístup a způsob napojení na
konkrétní rozhodnutí.

\subsection*{Shrnutí ke zkoušce}
\begin{itemize}
  \item Externí zdroje doplňují interní data o tržní, regionální, makroekonomický, regulatorní nebo
  konkurenční kontext.
  \item Primární zdroj přináší původní data, sekundární zdroj je shrnuje, interpretuje nebo katalogizuje.
  \item Důležité jsou oficiální statistiky, veřejné registry, národní a evropské datové portály,
  ČNB ARAD, RÚIAN a registry smluv nebo zakázek.
  \item Otevřená data mají být strojově čitelná, licenčně použitelná, dokumentovaná a dohledatelná
  v katalogu.
  \item Před použitím je nutné posoudit důvěryhodnost, metodiku, aktuálnost, granularitu, kvalitu
  a právní podmínky zdroje.
  \item Přístup může probíhat ručním stažením, hromadným exportem, API, databázovým přístupem,
  scrapingem nebo nákupem dat.
  \item Competitive intelligence je legální a etická práce s informacemi o vnějším prostředí
  organizace za účelem lepšího rozhodování.
  \item Proces competitive intelligence zahrnuje plánování, sběr, vyhodnocení, analýzu, distribuci
  a zpětnou vazbu.
  \item Externí data v pipeline vyžadují ukládání metadat, verzování, validace, monitoring a jasné
  mapování na interní data.
\end{itemize}

\subsection*{Zdroje}
\begin{itemize}
  \item Portál o datech České republiky: \href{https://data.gov.cz/pro-u\%C5\%BEivatele/otev\%C5\%99en\%C3\%A1-data/}{Otevřená data}.
  \item Český statistický úřad: \href{https://opendata.csu.gov.cz/}{Otevřená data ČSÚ}.
  \item Eurostat: \href{https://ec.europa.eu/eurostat/web/user-guides/data-browser/api-data-access/api-getting-started/api}{API - Getting started with statistics API}.
  \item Evropská komise: \href{https://digital-strategy.ec.europa.eu/en/factpages/open-data-and-high-value-datasets-step-step-access-guide}{Open data and high-value datasets}.
  \item Ministerstvo financí ČR: \href{https://mf.gov.cz/cs/ministerstvo/informacni-systemy/ares}{ARES - Administrativní registr ekonomických subjektů}.
  \item Ministerstvo spravedlnosti ČR: \href{https://verejnerejstriky.msp.gov.cz/clanky/obecne-informace/verejny-rejstrik-a-sbirka-listin}{Veřejný rejstřík a sbírka listin}.
  \item Česká národní banka: \href{https://www.cnb.cz/cs/statistika/arad-system-casovych-rad/}{ARAD - systém časových řad}.
  \item Česká národní banka: \href{https://www.cnb.cz/docs/arad20/dokumentace/arad_rest_api_cs.pdf}{ARAD-REST-API}.
  \item Intelligence.gov: \href{https://www.intelligence.gov/how-the-ic-works}{How the IC Works}.
  \item SAGE Encyclopedia of Business Ethics and Society: \href{https://sk.sagepub.com/ency/edvol/embed/sage-encyclopedia-of-business-ethics-and-society-2e/chpt/competitive-intelligence}{Competitive Intelligence}.
\end{itemize}

\newpage
